% Staged only. No canonical file has been edited.
% Requires amsmath, amssymb, amsthm, and the existing esn theorem labels.
% All new labels have namespace n24:.

\subsection{An exact cyclic-triality isometry of the entire Niemeier space}
\label{n24:sec:cyclic-isometry}

Retain the ordered fibres
\(f_i289^j\), \(1\leq i\leq4\), \(0\leq j\leq5\),
the Euclidean augmentation coordinates of \(A(F_{\eta_i})\), the
ordered orthonormal coordinates of \(D\subset\mathbb R^4\), and the
72-element glue \(\mathcal H\) of
\eqref{esn:eq:glue-and-overlattice}.  Thus the actual real span is
\[
 V_N=N\otimes_{\mathbb Z}\mathbb R
 =\left\{(x_{i,j};d_k)\in\mathbb R^{24}\oplus\mathbb R^4:
               \sum_{j=0}^5x_{i,j}=0\ (1\leq i\leq4)\right\},
\]
with squared norm \(\sum_{i,j}x_{i,j}^2+\sum_kd_k^2\).
The four displayed linear constraints, rather than an implicit change of
metric, give its dimension 24.  Choose an ordered orthonormal octonion
basis \((u_0,\ldots,u_7)\) with \(u_0=1\).  Every octonion product
below is the product of the chosen real octonion algebra, without a
change of multiplication convention.

For \(p\in\{0,1\}\) and \(j\in\{0,1,2\}\), set
\begin{align}
 m_{i,p}&=\frac13\sum_{j=0}^2x_{i,p+2j},&
 t_{i,p,j}&=x_{i,p+2j}-m_{i,p},\label{n24:eq:triplets}\\
 \alpha_i&=\frac{\sum_{j=0}^2x_{i,2j}
                    -\sum_{j=0}^2x_{i,1+2j}}{\sqrt6},&
 (b_0,\ldots,b_7)&=(\alpha_1,\alpha_2,\alpha_3,\alpha_4,
                           d_1,d_2,d_3,d_4).
 \label{n24:eq:fixed-coordinates}
\end{align}
In particular \(m_{i,0}=\alpha_i/\sqrt6\),
\(m_{i,1}=-\alpha_i/\sqrt6\), and
\(\sum_jt_{i,p,j}=0\).  The ordered map is
\begin{equation}
 \begin{split}
 \mathcal L:V_N&\longrightarrow\mathbb O^3,\\
 (x;d)&\longmapsto(q_1,q_2,q_3),\\
 q_{j+1}&=\sum_{i=1}^4\sum_{p=0}^1
       \left(t_{i,p,j}+\frac{b_{2(i-1)+p}}{\sqrt3}\right)
                                          u_{2(i-1)+p}.
 \end{split}
 \label{n24:eq:full-isometry}
\end{equation}

\begin{theorem}[Actual rank-24 map, inverse, pairing, and residue action]
\label{n24:thm:full-isometry}
The map \(\mathcal L\) is a real linear isometry onto
\(\mathbb O^3\) with its product norm.  Write
\(q_{j+1}=\sum_{a=0}^7q_{j,a}u_a\), and define
\begin{equation}
 b_a=\frac1{\sqrt3}\sum_{j=0}^2q_{j,a},\qquad
 t_{i,p,j}=q_{j,2(i-1)+p}
                    -\frac13\sum_{h=0}^2q_{h,2(i-1)+p}.
 \label{n24:eq:inverse-first}
\end{equation}
Its inverse is exactly
\begin{equation}
 x_{i,p+2j}=t_{i,p,j}+
                 \frac{(-1)^p b_{i-1}}{\sqrt6},
 \qquad d_k=b_{k+3}.
 \label{n24:eq:inverse-second}
\end{equation}
If \(r=\rho(289^2)=\rho(361)\), with the residue action
of Theorem~\ref{esn:fa:thm:residue-weyl}, and
\(c(q_1,q_2,q_3)=(q_3,q_1,q_2)\), then
\begin{equation}
 \mathcal Lr=c\mathcal L,\qquad r^3=c^3=I,
 \qquad rN=N.
 \label{n24:eq:cyclic-intertwining}
\end{equation}
Thus \(\mathcal L|_N\) embeds the original even unimodular lattice,
with every pairing and glue class retained, into the actual three
octonion sectors and intertwines an actual order-three lattice
isometry with cyclic triality.
\end{theorem}

\begin{proof}
The two means sum to zero because each six-coordinate fibre has total
sum zero.  The equations for \(m_{i,p}\) follow from
\(\alpha_i=6m_{i,0}/\sqrt6\).  Orthogonality to constant triplets gives
\[
 \sum_{j=0}^5x_{i,j}^2
 =\sum_{p=0}^1\sum_{j=0}^2t_{i,p,j}^2
                    +3m_{i,0}^2+3m_{i,1}^2
 =\sum_{p,j}t_{i,p,j}^2+\alpha_i^2.
\]
For each \(a=2(i-1)+p\), the same zero-sum identity gives
\[
 \sum_{j=0}^2\left(t_{i,p,j}+\frac{b_a}{\sqrt3}\right)^2
                      =\sum_{j=0}^2t_{i,p,j}^2+b_a^2.
\]
Summing over all eight values of \(a\) proves
\(\sum_{h=1}^3|q_h|^2=\sum_{i,j}x_{i,j}^2+\sum_kd_k^2\).
Polarization proves equality of all pairings.  Equations
\eqref{n24:eq:inverse-first} recover the means and residuals of each
output triplet.  Equation~\eqref{n24:eq:inverse-second} then recovers
the original means, the original fibre coordinates, and \(d\).
Conversely those formulas applied to arbitrary \(q\) yield
\(\sum_jx_{i,j}=0\) and return the prescribed \(q\).  This proves
bijectivity without a dimension-only argument.

The residue action sends the coordinate vector \(e_{f_i289^h}\)
to \(e_{f_i289^{h+2}}\), so the coefficient in output position
\(p+2j\) is the old coefficient in position \(p+2(j-1)\), with
\(j\) read modulo three.  Both means, all \(\alpha_i\), and all
\(d_k\) are fixed.  Hence the three output octonions become
\((q_3,q_1,q_2)\), proving the intertwining formula.  Each coordinate
permutation acts trivially on \(A^\vee/A\): for
\(\lambda=e_0-\frac16\sum_he_h\), its difference from \(\lambda\)
is an augmentation root.  The \(D\)-coordinates are fixed.  Thus
\(r\) acts trivially on the full discriminant group, preserves
\(\mathcal H\), and preserves its exact inverse image \(N\).
Its order is three, as seen on \(e_0-e_1\) in any fibre.
\end{proof}

\paragraph{The exact octonion-coordinate glue test.}
This embedding has the following fully explicit image, denoted
\(\Lambda_{\mathbb O}=\mathcal L(N)\).  Given \(q\), first recover
\((x;d)\) by \eqref{n24:eq:inverse-first}--\eqref{n24:eq:inverse-second}.
For each fibre require
\(x_{i,h}-x_{i,k}\in\mathbb Z\) for every \(h,k\), and put
\[
 k_i=-6x_{i,0}\pmod6\in\mathbb Z/6\mathbb Z.
\]
This is defined because
\(6x_{i,0}=\sum_h(x_{i,0}-x_{i,h})\in\mathbb Z\).
Require also that the four \(d_k\) are all integral or all
half-integral.  Let \(y(d)\in\mathbb F_2^2\) be the unique class with
\[
 d+D=0+D,\ v+D,\ s+D,\ c+D
 \quad\hbox{marked respectively by }00,10,01,11.
\]
Then the final and exact test is
\begin{equation}
 (k_1,k_2,k_3,k_4;y(d))
   =\bigl(3u+2(a+b,a+2b,a,b);\tau(u)\bigr)
 \label{n24:eq:image-glue}
\end{equation}
for \(u\in\mathbb F_2^4\) of even weight,
\((a,b)\in\mathbb F_3^2\), with the first four coordinates read
modulo six, and \(\tau(u)=(u_2+u_4,u_3+u_4)\).
Indeed the integral-difference conditions are precisely membership in
\(A^\vee\), and the common fractional part equals
\(-k_i/6\pmod{\mathbb Z}\), which is the fractional part of
\(k_i\lambda\).  The half-integral test is precisely membership in
\(D^\vee\), and \eqref{n24:eq:image-glue} is the original
\(\mathcal H\).  This proves both directions of the image
description.  No new lattice or rescaled glue is substituted.

\paragraph{The determinant and the factor two in the Albert metric.}
The matrix convention, also specifying this construction independently
of its host document, is
\[
 Q_{\mathbb O}(\lambda;q)=
 \begin{pmatrix}
 \lambda&q_3&\overline q_2\\
 \overline q_3&\lambda&q_1\\
 q_2&\overline q_1&\lambda
 \end{pmatrix}\in\mathfrak h_3(\mathbb O),\qquad
 X\circ Y=(XY+YX)/2.
\]
Put
\[
 E(q)=Q_{\mathbb O}(0;q_1,q_2,q_3),\qquad
 \mathfrak t(q_1,q_2,q_3)=\operatorname{Re}((q_1q_2)q_3),
 \qquad \mathfrak t_N(X)=\mathfrak t(\mathcal LX).
\]
The original ordered Albert determinant becomes the exact polynomial
\begin{equation}
 \det_{\mathbb O}\bigl(\lambda I+E(\mathcal LX)\bigr)
   =\lambda^3-\lambda(X,X)+2\mathfrak t_N(X).
 \label{n24:eq:niemeier-albert-cubic}
\end{equation}
Moreover
\begin{equation}
 \operatorname{tr}\bigl(E(\mathcal LX)\circ E(\mathcal LY)\bigr)
                  =2(X,Y).
 \label{n24:eq:albert-metric-two}
\end{equation}
For completeness, direct binary multiplication gives
$\operatorname{tr}E=0$, $\operatorname{tr}(E\circ E)=2\sum_i|q_i|^2$,
and $\operatorname{tr}(E\circ(E\circ E))=6\mathfrak t(q)$.
For the last equality, the three cyclic off-diagonal coordinates of
$E\circ E$ are $\overline{q_2q_3},\overline{q_3q_1},
\overline{q_1q_2}$; their trace pairings with the corresponding
$q_i$ each give $2\mathfrak t(q)$.
The Jordan trace formula
$\det X=\operatorname{tr}(X\circ(X\circ X))/3
-\operatorname{tr}X\operatorname{tr}(X\circ X)/2
+(\operatorname{tr}X)^3/6$
therefore gives the first identity after expanding $X=\lambda I+E$
and substituting the proved norm equality. For the second,
each off-diagonal octonion
occurs in both transposed matrix entries, giving exactly
\(2\sum_i\operatorname{Re}(q_i\overline{p_i})\).
Real associativity and cyclicity of octonion real parts give
\(\mathfrak t(q_3,q_1,q_2)=\mathfrak t(q_1,q_2,q_3)\):
\[
 \operatorname{Re}((q_3q_1)q_2)
 =\operatorname{Re}(q_3(q_1q_2))
 =\operatorname{Re}((q_1q_2)q_3).
\]
Consequently \(\mathfrak t_N(rX)=\mathfrak t_N(X)\), and the
whole determinant in \eqref{n24:eq:niemeier-albert-cubic} is
\(r\)-invariant with \(\lambda\) fixed.  Equivalently, for the real
permutation matrix \(P e_i=e_{i+1}\), indices modulo three,
\[
 P\bigl(\lambda I+E(q)\bigr)P^T
             =\lambda I+E(q_3,q_1,q_2).
\]
This is also a Jordan automorphism: multiplication by a real
permutation matrix merely reindexes the two-factor matrix products,
so it preserves \((AB+BA)/2\).  It therefore carries the positive
cone and its determinant boundary to themselves.  The cyclic bridge
is simultaneously linear, isometric on \(\mathbb O^3\), integral
on the stated lattice image, and compatible with the actual Albert
cubic and positive cone.

\paragraph{Dependence and the four-dimensional lattice summand.}
The map uses precisely the four specified origins \(f_i\), the
generator \(289\), the even/odd triplet order, the displayed order
of \(b\), the ordered \(D\)-coordinates, and the chosen orthonormal
octonion basis.  Changing one of these changes the displayed map;
the construction makes no claim of independence from these data.
Its restriction to the actual \(D\)-summand is
\begin{equation}
 d\longmapsto\frac1{\sqrt3}(z(d),z(d),z(d)),\qquad
 z(d)=d_1u_4+d_2u_5+d_3u_6+d_4u_7.
 \label{n24:eq:d-diagonal-embedding}
\end{equation}
This is an isometry on \(D\), and on \(D^\vee\) in its real span.
In particular its 24 shortest nonzero dual vectors become 24
specified points in this four-dimensional diagonal subspace.
The separate Cartan-weight map \(d\mapsto\sum_kd_k\varepsilon_k\)
is related to it by the exact isometry
\(\sum_kd_k\varepsilon_k\mapsto(z(d),z(d),z(d))/\sqrt3\).

\subsection{The Leech neighbour in the same three octonion coordinates}
\label{n24:sec:leech-coordinates}

The preceding map also transports the entire neighbour construction,
not only the real span of $N$. Put $N_{\mathbb O}=\mathcal L(N)$ and
retain the actual $\rho,\alpha,K,\Lambda$ of
Theorem~\ref{esnl:thm:leech-neighbor}. In the already fixed basis let
\[
 U=\sum_{a=0}^7u_a,\qquad
 V=\frac{u_0+u_1+u_2+u_3}{\sqrt2}
       +\sqrt3u_4+\frac{2u_5+u_6}{\sqrt3}.
\]
Substitution in the original isometry gives the explicit vectors
\begin{equation}
 \begin{aligned}
 \varrho=\mathcal L\rho&=(V+2U,V,V-2U),\\
 \delta=\mathcal L\alpha
 &=\left(\frac{(2+\sqrt2)u_0-2u_1}{3},
          \frac{(\sqrt2-1)u_0+u_1}{3},
          \frac{(\sqrt2-1)u_0+u_1}{3}\right).
 \end{aligned}
 \label{n24:eq:weyl-in-octonions}
\end{equation}
Indeed each even triplet of $\rho_A$ has mean $1/2$, each odd
triplet has mean $-1/2$, and both residual triplets are $(2,0,-2)$.
The four mean-difference coordinates are $\sqrt6/2$, while the
four $D$ coordinates are $(3,2,1,0)$. For $\alpha$, the two residual
triplets in the first fibre are $(2,-1,-1)/3$ and $(-2,1,1)/3$;
its only nonzero mean-difference coordinate is $2/\sqrt6$.
These facts prove every coefficient in
\eqref{n24:eq:weyl-in-octonions}. In particular
$|\varrho|^2=84$, $|\delta|^2=2$, and
$\langle\varrho,\delta\rangle=1$ in the product octonion metric.

Define
\begin{equation}
 \begin{aligned}
 K_{\mathbb O}&=\{q\in N_{\mathbb O}:
                          \langle\varrho,q\rangle\equiv0\pmod6\},\\
 \Lambda_{\mathrm{Leech},\mathbb O}
       &=K_{\mathbb O}+\mathbb Z(\varrho/6-\delta).
 \end{aligned}
 \label{n24:eq:leech-in-octonions}
\end{equation}
The isometry and its inverse prove, in both directions,
$\mathcal L(K)=K_{\mathbb O}$ and
$\mathcal L(\Lambda)=\Lambda_{\mathrm{Leech},\mathbb O}$.
Thus \eqref{n24:eq:leech-in-octonions} is a full coordinate
description of the Leech lattice inside the very same three sectors.
The glue test above, followed by its displayed height congruence,
is an exact membership test for $K_{\mathbb O}$; adjoining the six
cosets $k(\varrho/6-\delta)$, $0\le k<6$, gives all of the Leech
lattice. The common-sublattice indices, discriminant form, and
order-twelve quotient are transported without alteration.

The cyclic action also acts on the marked family of neighbours.
Retain $r=\rho(361)$ and its corresponding map
$c(q_1,q_2,q_3)=(q_3,q_1,q_2)$.
For any transported Weyl vector $\rho'=T\rho$, $T\in O(N)$,
the neighbour is defined using $T\alpha$, so that
$\langle\rho',T\alpha\rangle=1$. Then
\[
 c\,\mathcal L(K_{\rho'})=\mathcal L(K_{r\rho'}),\qquad
 c\,\mathcal L(\Lambda_{\rho'})=\mathcal L(\Lambda_{r\rho'}).
\]
This follows directly from $\mathcal Lr=c\mathcal L$ and
(NL18). The inverse maps use $r^{-1}$ and $c^{-1}$.
It specifies exactly which marked neighbour is the target; it does
not silently assume that the chosen $\rho$ is fixed by $r$.
Finally the common cubic evaluates on every one of these lattice
points as
\[
 \det Q_{\mathbb O}(\lambda;\mathcal LX)
    =\lambda^3-\lambda(X,X)
                   +2\operatorname{Re}(((\mathcal LX)_1
                                         (\mathcal LX)_2)(\mathcal LX)_3).
\]
This identity holds for all $X\in N_{\mathbb R}$, hence also
on the explicitly embedded $N$, $K$ and $\Lambda$.

\subsection{Exactly which root isometries preserve this glue}
\label{n24:sec:glue-symmetry}

Retain the coordinate convention \((P_px)_i=x_{p(i)}\).
An arbitrary root-lattice isometry has the original form
\[
 T(X_1,X_2,X_3,X_4;d)
   =\bigl((\varepsilon_iQ_{\sigma_i}X_{p(i)})_{i=1}^4;\ T_Dd\bigr),
\]
where \(\varepsilon_i\in\{\pm1\}\), \(\sigma_i\in S_6\),
\(p\in S_4\), and \(T_D\in O(D)\).  If \(B=d_D(T_D)\),
the necessary and sufficient condition for \(T(N)=N\) is
\begin{equation}
 B\tau(u)=\tau(P_pu)\ (u\in E),\qquad
 \bigl(\varepsilon_i\ell_{p(i)}\bigr)_{i=1}^4
                  =\bigl(\ell_iA\bigr)_{i=1}^4
 \quad\hbox{for some }A\in GL_2(\mathbb F_3),
 \label{n24:eq:all-glue-preservers}
\end{equation}
where
\(\ell_1=(1,1),\ell_2=(1,2),\ell_3=(1,0),\ell_4=(0,1)\).
The \(A\), and hence \(B=B_p\), are unique.  Indeed the
two-primary part of the glue is the graph of \(\tau\), and signs
are trivial modulo two.  Its three-primary part is the image of the
injective map \(j(a,b)=(a+b,a+2b,a,b)\), so preserving it is
equivalent to \(M_\varepsilon P_pj=jA\).  These are exactly
the two conditions in \eqref{n24:eq:all-glue-preservers}; the
primary decomposition proves their sufficiency as well as necessity.
Since the four \(\ell_i\) represent the four distinct projective
row lines over \(\mathbb F_3\), each \(A\) gives exactly one
signed permutation by this equation.  This recovers the precise
48-element glue quotient, not merely its action on three labels.

The canonical full-automorphism theorem proves, with these same
coordinates and section, that
\[
 O(N)=\bigl(S_6^4\times W(D_4)\bigr)\rtimes L(G),
 \qquad G\cong GL_2(\mathbb F_3).
\]
Conjugation by the actual map \(\mathcal L\) now gives the entire
orthogonal stabilizer of the embedded lattice, with explicit matrices
obtained from \eqref{n24:eq:full-isometry} and its inverse:
\begin{equation}
 O(\Lambda_{\mathbb O})
       =\mathcal L O(N)\mathcal L^{-1},\qquad
 \Gamma_G(g)=\mathcal LL_g\mathcal L^{-1},\qquad
 \mathcal LL_g=\Gamma_G(g)\mathcal L.
 \label{n24:eq:full-transported-action}
\end{equation}
Both inclusions in the first equality follow by conjugating an
isometry preserving the specified lattice.  This is a faithful
action on the whole 24-dimensional space. Each \(\Gamma_G(g)\)
commutes with \(c\),
because the canonical equation
\(L_g\rho(289^2)=\rho(289^2)L_g\) conjugates to that equation.

For completeness, the previously displayed \(D_4\)-triality
section and geometric frame permutations have a precise comparison
even though that particular proposed linear conjugacy fails.  On
\(V_N\), each of \(\widehat R,\widehat H\) has trace
\(5-5+2=2\): the transposed \(A_5\) pair contributes zero,
the other two components contribute \(5,-5\), and the
\(D_4\) reflection contributes 2.  Their three-cycle has trace
\(5+1=6\).  A geometric Albert-frame transposition acts on
\(\mathbb O^3\) by swapping two sectors with conjugation and by
conjugating the fixed sector; its trace is
\(1-7=-6\).  Its three-cycle has trace zero.  Therefore an
invertible map cannot conjugate these two particular \(S_3\)
representations, since trace is invariant under conjugation.
The complete real character decompositions are, respectively,
\[
 V_N|_{\langle\widehat R,\widehat H\rangle}
       \cong7\mathbf1\oplus5\operatorname{sgn}\oplus6\operatorname{Std},
 \qquad
 \mathbb O^3|_{S_3^{\rm frame}}
       \cong\mathbf1\oplus7\operatorname{sgn}\oplus8\operatorname{Std}.
\]
These multiplicities follow by taking inner products with the
characters \((1,1,1),(1,-1,1),(2,0,-1)\), on classes of sizes
\(1,3,2\).  The exact whole-space relation for the original section
is \eqref{n24:eq:full-transported-action}; the exact
determinant-compatible cyclic relation is
\eqref{n24:eq:cyclic-intertwining}.  In the latter relation the
residue element has trivial discriminant action while its octonion
action cycles the sectors.  Those are two explicitly calculated
actions of the same operator on two explicitly different constructions;
no identification between their quotient homomorphisms is assumed.

Finally, the connected group \(\operatorname{Spin}(8)\) acts on
the ambient triality module, and its subgroup preserving
\(\Lambda_{\mathbb O}\) is exactly its intersection with
\eqref{n24:eq:full-transported-action}.  It is finite: the columns
of a lattice isometry in a fixed lattice basis have prescribed
bounded lengths, so each column has finitely many possible lattice
values.  A connected group preserving a discrete full lattice
pointwise as a set acts trivially, since each orbit of a lattice
vector is connected and discrete and the lattice spans.  The
nontrivial connected triality action therefore does not preserve
the whole lattice.  The proved cyclic subgroup \(\langle c\rangle\)
does preserve both the lattice and its Albert cubic.  Equations
\eqref{n24:eq:all-glue-preservers} and
\eqref{n24:eq:full-transported-action} specify the full glue
stabilizer; they do not assert that every element of that finite
orthogonal group also preserves the cubic.
